\chapter{Network Security from the Low-Level Perspective}
\label{ch:network}

\epigraph{``The network is not a source of truth. The network is a source of bytes.''}{}

\learninggoal{
	\item Understand network packet structure at the byte level (Ethernet, IP, TCP/UDP headers).
	\item Identify low-level vulnerabilities in network protocol implementations.
	\item Explain packet injection, socket-level attacks, and kernel-bypass networking security.
	\item Apply the memory-safety mindset to network parsers and protocol stacks.
	\item Understand the security implications of kernel-bypass networking (DPDK, XDP).
}

\section{Network Packets at the Byte Level}

Network security discussions often focus on high-level protocols---TLS, HTTP, OAuth. This chapter examines what happens at the raw byte level, where the abstraction gaps between protocol specifications and implementations create vulnerabilities.

An Ethernet frame arriving at a network interface is a sequence of bytes:

\begin{lstlisting}[caption=Ethernet frame (simplified)]
  | Destination MAC (6) | Source MAC (6) | EtherType (2) | Payload | FCS (4) |
\end{lstlisting}

Within the payload, an IPv4 packet:

\begin{lstlisting}[caption=IPv4 header (20 bytes)]
  | Version (4b) | IHL (4b) | DSCP (6b) | ECN (2b) | Total Length (16b) |
  | Identification (16b)    | Flags (3b) | Fragment Offset (13b) |
  | TTL (8b)    | Protocol (8b) | Header Checksum (16b)        |
  | Source IP (32b)                                            |
  | Destination IP (32b)                                       |
\end{lstlisting}

Within the IP payload, a TCP segment:

\begin{lstlisting}[caption=TCP header (20 bytes)]
  | Source Port (16b)          | Destination Port (16b)        |
  | Sequence Number (32b)                                     |
  | Acknowledgment Number (32b)                               |
  | Data Offset (4b) | Reserved (6b) | Flags (6b) | Window (16b) |
  | Checksum (16b)             | Urgent Pointer (16b)          |
\end{lstlisting}

These are the data structures that network parsers must process---and every parsing operation is a potential memory-safety bug.

\section{Protocol Parsing Vulnerabilities}

Every packet arriving at a host is parsed by the kernel's network stack (or a kernel-bypass stack). Parsing is where memory-safety bugs live:

\subsection{Length Field Validation}

Protocol headers contain length fields that must be validated before use:

\begin{lstlisting}[language=C, caption=Vulnerable IP length parsing]
int parse_ipv4(struct sk_buff *skb) {
    struct iphdr *ip = ip_hdr(skb);
    unsigned int hdr_len = ip->ihl * 4;  // IHL is 4-bit header length
    unsigned int tot_len = ntohs(ip->tot_len);
    // BUG: if hdr_len < 20 or tot_len < hdr_len:
    //   negative length, integer underflow, etc.
    if (skb->len < tot_len) return -EINVAL;   // Good check
    // But what if ip->ihl < 5 (minimum)? Then hdr_len < 20.
    // What if ip->ihl > 15? Then hdr_len > 60 (max header).
    parse_options(skb, ip, hdr_len);  // OOB read
}
\end{lstlisting}

The Linux kernel's IP parser validates these fields correctly in the mainline, but hundreds of embedded TCP/IP stacks (in IoT devices, network cards, boot ROMs) may not. Off-by-one errors in length validation are a common source of CVEs.

\subsection{TTL and Fragment Handling}

IP fragmentation introduces state into what would otherwise be a stateless parser. The kernel must reassemble fragments, which requires:

\begin{itemize}
	\item Tracking which fragments have arrived.
	\item Enforcing a maximum reassembly buffer size.
	\item Detecting overlapping fragments (which may be used to bypass firewalls).
	\item Timing out incomplete fragments.
\end{itemize}

Every aspect of this state machine has been vulnerable at some point. A classic attack is the \emph{teardrop} attack: overlapping IP fragments cause a calculation error in fragment length, leading to a kernel panic.

\section{Buffer Overflow in Network Parsers}

Network parsers are buffer-overflow hotspots because they must handle attacker-controlled lengths:

\begin{lstlisting}[language=C, caption=Network buffer overflow pattern]
int parse_protocol(const u8 *data, size_t len) {
    u8 option_buf[64];
    u16 opt_len = *(u16 *)(data + 2);  // Attacker controls this!
    if (opt_len > 64) return -EINVAL;   // Check? Maybe not.
    memcpy(option_buf, data + 4, opt_len);  // Overflow!
}
\end{lstlisting}

The Linux kernel handles tens of thousands of such length checks. Any one that is missing or incorrect is a potentially critical vulnerability.

\section{Socket-Level Attacks}

\subsection{TCP Desynchronisation}

An attacker who can observe or manipulate TCP traffic can desynchronise the TCP state machine by injecting packets with crafted sequence numbers. This can:

\begin{itemize}
	\item Cause data corruption by injecting forged segments.
	\item Reset an established connection (RST attack).
	\item Blindly inject data into a TCP stream (if the sequence number is guessed or observed).
\end{itemize}

\subsection{TCP Side Channels}

The TCP header's timestamp option (TSval/TSecr) can be used as a side channel to count packets or measure clock skew---useful for fingerprinting the sender's operating system or network topology.

\section{Kernel-Bypass Networking}

Traditional networking goes through the kernel's network stack, which provides safety (validation, isolation, scheduling) at the cost of performance. Kernel-bypass technologies (DPDK, XDP, AF\_XDP) give user-space programs direct access to the network interface card (NIC).

\begin{definition}[Kernel-Bypass]
	\textbf{Kernel-bypass networking} allows a user-space application to send and receive raw network frames without traversing the kernel's network stack. The application communicates directly with the NIC hardware.
\end{definition}

\subsection{Security Implications}

Kernel bypass trades safety for performance:

\begin{longtable}{p{3.5cm}p{5cm}p{4cm}}
	\toprule
	\textbf{Aspect}      & \textbf{Kernel networking} & \textbf{Kernel bypass}           \\
	\midrule
	Packet validation    & Kernel enforces checks     & Application's responsibility     \\
	Privilege separation & Kernel provides isolation  & Application has direct HW access \\
	Memory safety        & Kernel validator has CFI   & Application must validate        \\
	Mitigation access    & ASLR, NX, SMEP/SMAP, KPTI  & Same (kernel runs mitigations)   \\
	Bug impact           & Kernel memory corruption   & Still kernel context (via DMA)   \\
	\bottomrule
\end{longtable}

A bug in a DPDK packet parser is particularly dangerous because:

\begin{enumerate}
	\item The parser runs in user space but has access to DMA buffers---a bug gives direct hardware access.
	\item There is no kernel validator between the NIC and the application to catch errors.
	\item DPDK applications are typically written in C with extensive use of zero-copy (direct buffer manipulation), which amplifies the impact of memory-safety bugs.
\end{enumerate}

\subsection{eXpress Data Path (XDP)}

XDP attaches a BPF program to the network interface, running before the kernel's network stack. This allows:

\begin{itemize}
	\item Early packet filtering (faster than iptables/nftables).
	\item Packet modification (header rewrite, encapsulation).
	\item Redirection to user space via AF\_XDP sockets.
\end{itemize}

XDP programs run in a BPF sandbox with verifier-enforced safety: bounded loops, no arbitrary memory writes, and termination guarantees. This makes XDP safer than DPDK for the same use case.

\section{Wireless and Low-Level Attacks}

Wireless networks add another layer of low-level complexity:

\begin{itemize}
	\item \textbf{Frame injection:} crafting raw 802.11 frames to perform deauthentication attacks, beacon flooding, or evil-twin access points.
	\item \textbf{Wi-Fi parser bugs:} the 802.11 frame format includes numerous optional fields (HT/VHT/HE capabilities, information elements)---each is a potential parsing vulnerability in the firmware or kernel driver.
	\item \textbf{Firmware-level exploits:} Wi-Fi chipsets run proprietary firmware that processes raw frames. Bugs in this firmware can give an attacker arbitrary code execution on the Wi-Fi chip itself, often with DMA access to host memory.
\end{itemize}

\begin{example}[Broadpwn (CVE-2017-9417)]
	A buffer overflow in Broadcom Wi-Fi firmware's handling of beacon frames allowed an attacker in range to execute code on the Wi-Fi chipset, which had full DMA access to the host's memory. No user interaction was required---only a nearby malicious access point.
\end{example}

\section{Exercises}

\begin{exercise}
	Given a raw Ethernet frame (as hex bytes), identify the source and destination MAC addresses, EtherType, and the IP version. Use a packet capture tool (tcpdump or Wireshark) to verify.
\end{exercise}

\begin{exercise}
	An embedded TCP/IP stack parses an incoming TCP segment. The header length field (Data Offset) is 15, implying a 60-byte header. The total segment length is 40 bytes. What happens when the parser tries to read header options? What vulnerability class is this?
\end{exercise}

\begin{exercise}
	Compare the security properties of XDP and DPDK for a packet-processing application. Under what circumstances would each be more secure?
\end{exercise}

\begin{exercise}
	Research CVE-2020-11989 (a Linux kernel vulnerability in the netfilter subsystem). Describe the bug, its root cause, and how it might be exploited.
\end{exercise}
